\documentclass[11pt]{article}
\usepackage[margin=1in]{geometry}
\usepackage{graphicx, booktabs, amsmath, amssymb, amsthm, hyperref, caption, longtable}
\usepackage{harvard}
\usepackage{tabularx}
\usepackage{booktabs}
\usepackage{array}
\usepackage{siunitx}
\usepackage{makecell}
\usepackage{graphicx}
\graphicspath{{figs/}{./}{/mnt/data\_analysista/}}
\sisetup{detect-weight=true,detect-family=true}
\citationstyle{dcu}
\renewcommand{\harvardurl}[1]{\texttt{#1}}

\title{CC-Bench Trajectory Analysis: Metric Definitions, Scoring Framework, and Model-wise Evaluation}
\author{}


\newtheorem{definition}{Definition}
\newcommand{\density}{\mathrm{density}}
\newcommand{\densityw}{\mathrm{density}_{w}}
\newcommand{\sentset}{S}
\newcommand{\plan}{\Pi_{\text{plan}}}
\newcommand{\verify}{\Pi_{\text{verify}}}
\newcommand{\fix}{\Pi_{\text{fix}}}

\begin{document}
\maketitle
\section{Overview}
We address the CC-Bench prompt directly: select representative tasks, compare two models on
their behavioral and reasoning trajectories, and justify each score through observed interaction
patterns. Specifically, we evaluate \textbf{Claude-Sonnet-4} and \textbf{DeepSeek-V3.1-Terminus}
on five data-analysis trajectories (IDs~40–44), reporting three behavioral indicators—
\textit{Goal}, \textit{Strategy}, and \textit{Bugfix}—alongside auxiliary metrics
(SDS, TCR, SCR, RD) and two composite measures: the \textit{E~score} and
\textit{$\Delta M$ index}. Each captures a distinct reasoning dimension:
goal framing, structural planning, and verification or correction. 
Weighted variants (Goal$_w$, Strategy$_w$, Bugfix$_w$, RD$_w$) apply
\textbf{time-scale weighting} to highlight sustained and late-phase reasoning.

The composite $\Delta M$ index integrates structural (SDS), verification (SCR), coherence (E),
and brevity (TCR) differences between models. For data-analysis tasks, verification-related
components (SDS and SCR) receive higher weights (0.35 each), coherence (E) a moderate weight (0.20),
and brevity (TCR) a smaller penalty (0.10), reflecting reliability and diagnostic stability
over linguistic compression. Each task pair is then labeled as \textbf{Good}, \textbf{Same}, or
\textbf{Bad} based on $\Delta M$, forming the GSB–w classification.

In addition, three \textbf{visualization indicators}—Goal Priming (G), Strategy Density (S),
and Bug/Verification Density (B)—facilitate intuitive inspection of model behavior.
These qualitative indicators are excluded from scoring but illustrate goal setting,
planning density, and verification strength for human interpretation and reporting.

Overall, reliability dominates communicative efficiency in these data-analysis trajectories.
Claude acts as a concise planner with higher goal clarity and compression (TCR),
while DeepSeek functions as a reflective verifier, maintaining denser verification
(Bugfix$_w$, SCR$_w$) and stronger late-phase reflection (RD$_w$).
Representative cases such as Task~44 show DeepSeek’s sustained self-checking yielding higher
E~score and a positive $\Delta M$ under reliability-oriented weighting.
All quantitative values are reported in Tables~1–3; metric derivations appear in Appendix~A.


\section{Results and Analysis}
\subsection{Quantitative Results}
Tables~\ref{tab:da-behavior} and~\ref{tab:da-efficiency} summarize the quantitative metrics across five representative \texttt{da} trajectories for two models.
Table~\ref{tab:da-behavior} reports the behavioral indicators (G, S, B, SDS, TCR, SCR, RD),
while Table~\ref{tab:da-efficiency} extends the view with efficiency and outcome indicators
(\textit{assistant\_text\_len}, \textit{E\_score}, \textit{delta\_M}, and the final \textit{GSB} verdict).
Together, they form the empirical basis for model-wise comparison and interaction-level analysis.

\vspace{6pt}

% --------- Table 2A ----------
\begin{table}[htbp]
\renewcommand{\arraystretch}{1.18}
\setlength{\tabcolsep}{5pt}
\centering
\small
\begin{tabularx}{\textwidth}{
  >{\raggedright\arraybackslash}p{0.8cm}
  >{\raggedright\arraybackslash}p{3.0cm}
  >{\raggedright\arraybackslash}p{1.0cm}
  *{7}{>{\centering\arraybackslash}X}
}
\toprule
\textbf{ID} & \textbf{Model} & \textbf{Task} &
\textbf{Goal} & \textbf{Strategy} & \textbf{Bugfix} & \textbf{SDS} & \textbf{TCR} & \textbf{SCR} & \textbf{RD} \\
\midrule
40 & DeepSeek{-}V3.1{-}Terminus & da & 0.300 & 0.001 & 0.031 & 0.166 & 0.458 & 0.035 & 0.000 \\
40 & Claude{-}Sonnet{-}4        & da & 0.300 & 0.000 & 0.009 & 0.110 & 0.644 & 0.015 & 0.000 \\
41 & DeepSeek{-}V3.1{-}Terminus & da & 0.300 & 0.000 & 0.075 & 0.101 & 0.421 & 0.076 & 0.002 \\
41 & Claude{-}Sonnet{-}4        & da & 0.300 & 0.000 & 0.038 & 0.067 & 0.678 & 0.038 & 0.000 \\
42 & DeepSeek{-}V3.1{-}Terminus & da & 0.300 & 0.000 & 0.027 & 0.104 & 0.595 & 0.033 & 0.000 \\
42 & Claude{-}Sonnet{-}4        & da & 0.300 & 0.000 & 0.067 & 0.138 & 0.613 & 0.069 & 0.000 \\
43 & DeepSeek{-}V3.1{-}Terminus & da & 0.300 & 0.000 & 0.020 & 0.063 & 0.337 & 0.020 & 0.000 \\
43 & Claude{-}Sonnet{-}4        & da & 0.300 & 0.001 & 0.059 & 0.086 & 0.586 & 0.059 & 0.000 \\
44 & DeepSeek{-}V3.1{-}Terminus & da & 0.300 & 0.000 & 0.081 & 0.150 & 0.544 & 0.087 & 0.000 \\
44 & Claude{-}Sonnet{-}4        & da & 0.347 & 0.005 & 0.038 & 0.114 & 0.661 & 0.042 & 0.000 \\
\bottomrule
\end{tabularx}
\vspace{4pt}
\caption{Table 2A: Behavioral metrics for five \texttt{da} trajectories. Abbreviations: G (Goal Priming), S (Strategy Density), B (Bug/Verification Density), SDS (Strategic Density Score), TCR (Textual Compression Ratio), SCR (Sanity/Correction Ratio), RD (Reflective Density).}
\label{tab:da-behavior}
\end{table}

\vspace{6pt}
% --------- Table 2B ----------
\begin{table}[htbp]
\renewcommand{\arraystretch}{1.18}
\setlength{\tabcolsep}{5pt}
\centering
\small
\begin{tabularx}{\textwidth}{
  >{\raggedright\arraybackslash}p{0.8cm}
  >{\raggedright\arraybackslash}p{3.0cm}
  >{\raggedright\arraybackslash}p{1.0cm}
  *{3}{>{\centering\arraybackslash}X}
  >{\raggedright\arraybackslash}p{0.9cm}
}
\toprule
\textbf{ID} & \textbf{Model} & \textbf{Task} &
\textbf{Len} & \textbf{E} & $\boldsymbol{\Delta M}$ & \textbf{GSB} \\
\midrule
40 & DeepSeek{-}V3.1{-}Terminus & da & 35072 & 0.0039849 & +0.045796 & Good \\
40 & Claude{-}Sonnet{-}4        & da & 19406 & 0.0010036 & -0.045796 & Bad  \\
41 & DeepSeek{-}V3.1{-}Terminus & da & 20752 & 0.0055440 & +0.051705 & Good \\
41 & Claude{-}Sonnet{-}4        & da & 15700 & 0.0015173 & -0.051705 & Bad  \\
42 & DeepSeek{-}V3.1{-}Terminus & da & 21360 & 0.0021517 & -0.023450 & Bad  \\
42 & Claude{-}Sonnet{-}4        & da & 22085 & 0.0059033 & +0.023450 & Good \\
43 & DeepSeek{-}V3.1{-}Terminus & da & 20592 & 0.0009424 & +0.002749 & Same \\
43 & Claude{-}Sonnet{-}4        & da & 16291 & 0.0031992 & -0.002749 & Same \\
44 & DeepSeek{-}V3.1{-}Terminus & da & 27933 & 0.0084521 & +0.041164 & Good \\
44 & Claude{-}Sonnet{-}4        & da & 14901 & 0.0028826 & -0.041164 & Bad  \\
\bottomrule
\end{tabularx}
\vspace{4pt}
\caption{Table 2B: Efficiency and verdict indicators. Abbreviations: Len (assistant\_text\_len), E (E\_score), $\Delta M$ (delta\_M).}
\label{tab:da-efficiency}
\end{table}

\subsection{Model-wise Behavioral Comparison}
Using Tables~\ref{tab:da-behavior} and~\ref{tab:da-efficiency}, we observe complementary reasoning tendencies across tasks 40--44:
\begin{itemize}
  \item \textbf{Overall verdict (GSB):} DeepSeek wins 3 tasks (40, 41, 44), Claude leads on one (42), and task 43 is tied — giving a \textbf{3--1--1} advantage to DeepSeek.
  \item \textbf{Verification and reflexivity:} DeepSeek’s advantage stems from higher $B$ and $SCR$ in most cases (e.g., 0.075/0.081 and 0.076/0.087), correlating with positive $\Delta M$ and stronger $E$-scores.
  \item \textbf{Claude’s Task 42 exception:} Claude outperforms with higher $B{=}0.067$ and $SCR{=}0.069$ and achieves the best $E{=}0.0059$, reflecting transient verification efficiency.
  \item \textbf{Compression vs.\ density trade-off:} Claude maintains higher \emph{TCR}, showing compact articulation, while DeepSeek achieves higher verification density—longer but more diagnostic reasoning.
  \item \textbf{Goal priming:} Both stay near $G\!\approx\!0.30$; Claude’s spike to 0.347 (Task 44) suggests sharper goal setup but insufficient follow-through in verification.
\end{itemize}
In sum, \textbf{Claude} functions as a \emph{concise planner}, whereas \textbf{DeepSeek} acts as a \emph{reflective verifier}. 
These dual archetypes underscore a balance between communicative economy and procedural depth, with DeepSeek favored in reliability-oriented settings.

\subsection{Interaction-Level Notes}
The efficiency metrics in Table~\ref{tab:da-efficiency} clarify trade-offs between brevity and stability.
DeepSeek produces longer assistant outputs (e.g., tasks 40, 44) yet converts that budget into higher $SCR$ and $E$.
Claude’s shorter turns (high TCR) deliver compactness but less verification gain, explaining its lower average $\Delta M$ and mixed GSB results.

\subsection{Visualization}
Figures~\ref{fig:combined-vis} and~\ref{fig:emergence-trend} visualize the updated, \textbf{weighted} behavioral and reasoning dynamics across the five \texttt{da} tasks.

The GSB chart (top of Figure~\ref{fig:combined-vis}) shows that both models maintain similar early goal clarity ($G{\approx}0.30$),
but \textbf{DeepSeek consistently exhibits higher verification density ($B$)} and a slightly more stable strategic pattern ($S$).
In contrast, \textbf{Claude displays a sharper goal-priming peak} at task~44 ($G{=}0.347$),
yet this early planning advantage is offset by weaker verification behavior, resulting in lower overall robustness.

The auxiliary metrics (bottom of Figure~\ref{fig:combined-vis}) highlight a clearer trade-off once time-weighted variants are considered.
\textbf{Claude maintains higher textual compression ($TCR$)} across all tasks—indicating concise articulation and tighter phrasing—
whereas \textbf{DeepSeek achieves higher weighted strategic density ($SDS_{w}$)} and verification ratio ($SCR_{w}$),
signaling more persistent reasoning and stronger self-correction throughout the dialogue trajectory.
Although both models show near-zero unweighted reflective density ($RD{\approx}0$),
the \textbf{time-weighted reflective density ($RD_{w}$)} slightly increases for DeepSeek,
suggesting that its reflective or corrective statements tend to appear in later conversational phases.
This temporal behavior aligns with DeepSeek’s higher $SCR_{w}$ and the delayed but stronger emergence peaks observed in Figure~\ref{fig:emergence-trend}.

Figure~\ref{fig:emergence-trend} plots the \textbf{weighted emergent reasoning scores} ($E\_score$) across the same task set.
A mild cross-over dynamic is observed: \textbf{Claude rises early at task~42} ($E{=}0.0055$),
while \textbf{DeepSeek peaks later at task~44} ($E{=}0.0074$), maintaining higher values across most earlier tasks.
This contrast reflects two complementary reasoning modes:
Claude’s compact, early-emerging reasoning efficiency versus DeepSeek’s delayed yet more iterative emergence,
driven by verification and late-phase reflection.
Together, these visual patterns support the quantitative findings from Tables~\ref{tab:da-behavior} and~\ref{tab:emergent-reasoning},
clarifying the balance between \emph{compression-oriented} and \emph{verification-oriented} reasoning architectures.


Figure~\ref{fig:emergence-trend} visualizes the emergent reasoning scores ($E\_score$) across tasks.
The updated curve reveals a \textbf{cross-over dynamic}: Claude peaks early at task~42 ($E{=}0.0059$), while DeepSeek climbs sharply to a higher peak at task~44 ($E{=}0.0085$).
This pattern illustrates their complementary reasoning modes:
Claude’s reasoning strength emerges through compactness and early planning,
whereas DeepSeek’s emergence is delayed but more pronounced—driven by iterative verification and reflective correction.
Together, these visual patterns corroborate the quantitative findings from Tables~\ref{tab:da-behavior} and~\ref{tab:da-efficiency},
clarifying the trade-off between \emph{compression-oriented} and \emph{verification-oriented} reasoning architectures.

\begin{figure}[htbp]
  \centering
  \includegraphics[width=.7\linewidth]{figs/gsb_by_model_unweighted.png}\\[1em]
  \includegraphics[width=.7\linewidth]{figs/aux_by_model_unweighted.png}
  \caption{Model-wise metric visualizations: (Top) GSB scores; (Bottom) auxiliary metrics (SDS/TCR/SCR/RD).}
  \label{fig:combined-vis}
\end{figure}

\clearpage
\section{Emergent Reasoning Capability Analysis}
This section examines \textbf{emergent reasoning} by linking structural indicators (SDS, TCR, SCR, RD) with efficiency and robustness (E\_score, active\_ratio, failure\_rate).
All data\_analysista are drawn from representative \texttt{data\_analysis} tasks (IDs 40--44).

\subsection{Quantitative Summary}

\begin{table}[htbp]
\renewcommand{\arraystretch}{1.22}
\centering
\small
\begin{tabularx}{\textwidth}{
  >{\raggedright\arraybackslash}p{3.2cm}
  >{\centering\arraybackslash}p{0.9cm}
  S[table-format=1.3]
  S[table-format=1.3]
  S[table-format=1.3]
  S[table-format=1.3]
  S[table-format=1.5]
  S[table-format=1.3]
  S[table-format=2.2]
  S[table-format=2.1]
}
\toprule
\textbf{Model} & \textbf{Task} &
\multicolumn{5}{c}{\textbf{Behavioral (Weighted Variants)}} &
\multicolumn{3}{c}{\textbf{Efficiency / Robustness}}\\
\cmidrule(lr){3-7}\cmidrule(lr){8-10}
& &
\multicolumn{1}{c}{\textbf{SDS$_w$}} &
\multicolumn{1}{c}{\textbf{TCR}} &
\multicolumn{1}{c}{\textbf{SCR$_w$}} &
\multicolumn{1}{c}{\textbf{RD$_w$}} &
\multicolumn{1}{c}{\textbf{E\_score}} &
\multicolumn{1}{c}{\makecell{\textbf{active}\\\textbf{ratio}}} &
\multicolumn{1}{c}{\makecell{\textbf{avg\_turn}\\\textbf{time (s)}}} &
\multicolumn{1}{c}{\makecell{\textbf{failure}\\\textbf{rate (\%)}}} \\
\midrule
DeepSeek-V3.1-Terminus & 40 & 0.183 & 0.458 & 0.048 & 0.000 & 0.00602 & 0.823 & 22.34 & 15.8 \\
Claude-Sonnet-4        & 40 & 0.089 & 0.644 & 0.012 & 0.000 & 0.00065 & 0.763 & 68.01 &  7.1 \\
DeepSeek-V3.1-Terminus & 41 & 0.100 & 0.421 & 0.085 & 0.007 & 0.00647 & 0.794 & 15.79 &  8.3 \\
Claude-Sonnet-4        & 41 & 0.081 & 0.678 & 0.061 & 0.000 & 0.00294 & 0.340 & 11.67 & 18.8 \\
DeepSeek-V3.1-Terminus & 42 & 0.077 & 0.595 & 0.038 & 0.000 & 0.00183 & 0.774 & 15.39 & 19.2 \\
Claude-Sonnet-4        & 42 & 0.139 & 0.613 & 0.064 & 0.000 & 0.00552 & 0.490 & 25.85 & 24.1 \\
DeepSeek-V3.1-Terminus & 43 & 0.097 & 0.337 & 0.039 & 0.000 & 0.00283 & 0.782 & 14.15 &  9.4 \\
Claude-Sonnet-4        & 43 & 0.052 & 0.586 & 0.019 & 0.000 & 0.00062 & 0.676 & 40.27 & 12.9 \\
DeepSeek-V3.1-Terminus & 44 & 0.145 & 0.544 & 0.079 & 0.000 & 0.00742 & 0.685 & 13.27 & 10.9 \\
Claude-Sonnet-4        & 44 & 0.185 & 0.661 & 0.076 & 0.000 & 0.00846 & 0.322 &  9.05 & 22.2 \\
\bottomrule
\end{tabularx}
\caption{Emergent reasoning indicators (weighted variants) and efficiency statistics across five \texttt{da} tasks (IDs 40–44). Weighted forms (SDS$_w$, SCR$_w$, RD$_w$) emphasize reasoning persistence and late-stage reflection.}
\label{tab:emergent-reasoning}
\end{table}


\subsection{Comparative Analysis}
Several cross-metric tendencies stand out:
\begin{itemize}
  \item \textbf{Weighted strategic density:} DeepSeek maintains higher $SDS_w$ in 40, 41, 44 (0.10–0.18), showing richer structured reasoning across longer trajectories.
  \item \textbf{Verification and reflection:} DeepSeek’s $SCR_w$ rises to 0.048–0.085 and includes non-zero $RD_w$ at task 41 (0.007), whereas Claude’s reflective behavior remains near zero, confirming deeper self-checking in DeepSeek.
  \item \textbf{Claude’s compact emergence:} Claude attains higher $TCR$ and an early $E\_score$ peak at 42 (0.0055), emphasizing compression-oriented efficiency.
  \item \textbf{Temporal efficiency:} DeepSeek sustains higher active ratios (0.68–0.82) and shorter average turns (13–22 s), balancing reasoning depth and flow stability.
  \item \textbf{Failure robustness:} Both remain comparable, with DeepSeek more consistent (8–19 \%) and Claude more volatile (7–24 \%).
\end{itemize}


\subsection{Interpretation}
Overall, we observe a trade-off between \emph{compression-oriented reasoning} (Claude)
and \emph{verification-oriented reasoning} (DeepSeek).
DeepSeek’s higher SDS and SCR indicate stronger coherence over extended trajectories,
hinting at early-stage \textbf{emergent meta-reasoning}.
Claude, by contrast, favors textual compactness (high TCR) at the cost of iterative verification.
These differences align with distinct \emph{reasoning archetypes}---\emph{linear planners} versus \emph{reflective validata\_analysistors}.

\begin{figure}[htbp]
  \centering
  \includegraphics[width=.75\linewidth]{figs/emergence_trend.png}
  \caption{Reasoning emergence by task.
  Each point represents a data\_analysista-analysis task (IDs 40--44).
  The \textbf{E\_score} quantifies emergent reasoning strength from structural and temporal features
  (SDS, TCR, SCR, RD, and activity ratios).
  DeepSeek-V3.1-Terminus exhibits more volatile but higher peaks (notably Task~44),
  while Claude-Sonnet-4 shows steadier yet lower emergence across tasks.}
  \label{fig:emergence-trend}
\end{figure}

\section{Interpretation and Discussion}
Across the representative data\_analysista-analysis trajectories (Tasks 40--44), the two models reveal distinct reasoning styles under the GSB framework.
\textbf{Claude-Sonnet-4} shows slightly higher G and stronger text compression, reflecting an early and concise articulation of objectives followed by direct, linear execution.
Its responses are compact and plan-first, efficient in setting up intentions but less persistent in iterative reasoning.

In contrast, \textbf{DeepSeek} demonstrates higher B and moderately stronger S across multiple tasks.
This pattern aligns with its trajectory: DeepSeek typically begins with a clear goal ($G \approx 0.30$) and then engages in self-checking and correctional steps --- using phrases like ``check'', ``verify'', or ``fix'' --- showing a reflective, verification-oriented process.
The higher B (0.075--0.081 vs.\ 0.038--0.059) captures this behavior quantitatively.

\textbf{Evaluation logic.} Within GSB, G reflects goal clarity, S reflects structured reasoning, and B captures self-monitoring and correction.
Since these tasks emphasize reliability in data\_analysista analysis, B carries greater interpretive weight.
Although Claude excels in brevity and clarity of initial planning, DeepSeek’s repeated self-verification and reflective adjustments contribute more to reasoning stability.

\textbf{Overall interpretation.} The two models embody complementary cognitive styles:
\begin{itemize}
  \item \textbf{Claude-Sonnet-4 $\rightarrow$ Concise Planner:} favors linear, efficient execution with minimal redundata\_analysisncy.
  \item \textbf{DeepSeek $\rightarrow$ Reflective Verifier:} adopts an iterative validata\_analysistion loop, maintaining steadier reasoning over longer trajectories.
\end{itemize}

In scenarios requiring procedural correctness and diagnostic depth, \textbf{DeepSeek achieves a higher overall GSB score}, whereas \textbf{Claude} emphasizes communicative economy.
This contrast highlights a broader trade-off between \emph{compression-oriented} and \emph{verification-oriented} reasoning that we see in many modern assistants.



\clearpage
\appendix
\section{Metric Framework and Interpretive Logic}
\label{appendix:metric-framework}

This appendix presents the metric definitions and the interpretive logic used throughout the paper.
It formalizes how each behavioral indicator (G, S, B, SDS, TCR, SCR, RD) is computed and how these
metrics map to cognitive and reasoning behaviors.

\subsection{Core Lexicons}
We define three regex lexicons that capture planning, verification, and debugging intents:
\begin{align*}
\plan&:\ \texttt{plan|step|strategy|outline|pipeline|procedure|first|then|next|we will} \\
\verify&:\ \texttt{check|verify|test|assert|unit test|print|plot|compare|evaluate} \\
\fix&:\ \texttt{fix|bug|error|patch|adjust|correct|refactor|resolve|rerun}
\end{align*}

\subsection{Core Metrics (GSB)}
\begin{definition}[Goal Priming ($G$)]
Measures whether the assistant initiates planning early:
\[
G = \min\!\bigl(1,\ 0.3 + 0.7 \cdot \density(\texttt{first\_chunk}, \plan)\bigr).
\]
\end{definition}

\begin{definition}[Strategy Density ($S$)]
\[
S = \density(\texttt{assistant\_text}, \plan),
\]
representing sustained planning and structured reasoning.
\end{definition}

\begin{definition}[Bug/Verification Density ($B$)]
\[
B = \max\!\bigl(\density(\texttt{assistant\_text}, \verify),\ \density(\texttt{assistant\_text}, \fix)\bigr),
\]
capturing self-checking and debugging tendencies.
\end{definition}

\paragraph{Weighted variants.}
Weighted forms $(G_w, S_w, B_w, RD_w)$ use $\densityw$, where time-scale weights
emphasize behaviors occurring later in the assistant trajectory. This adjustment is especially relevant for reflective behaviors ($RD_w$), which tend to emerge in later phases.

\subsection{Auxiliary Metrics}
\begin{itemize}
\item \textbf{SDS (Strategic Density Score):} Fraction of planning, verification, and fixing sentences within assistant outputs.
\item \textbf{TCR (Textual Compression Ratio):} Ratio of unique to total sentences (higher means less redundancy and greater conciseness).
\item \textbf{SCR (Sanity/Correction Ratio):} $\density(\texttt{assistant\_text}, \verify \cup \fix)$, measuring the frequency of validation or correction acts.
\item \textbf{RD (Reflective Density):} Proportion of reflective cues (``wait'', ``rethink'', ``revisit'', ``oops''). A time-weighted variant $RD_w$ captures late-phase reflection.
\end{itemize}

\subsection{Cognitive and Behavioral Interpretation}
Table~\ref{tab:metric-interpretation} summarizes how each metric family maps to reasoning abilities and observable behavior.

\begin{table}[htbp]
\renewcommand{\arraystretch}{1.3}
\setlength{\extrarowheight}{2pt}
\centering
\scriptsize
\begin{tabularx}{\textwidth}{
  >{\raggedright\arraybackslash}p{2.2cm}
  >{\raggedright\arraybackslash}p{3.2cm}
  >{\raggedright\arraybackslash}p{3.2cm}
  >{\raggedright\arraybackslash}X
}
\toprule
\textbf{Metric / Group} & \textbf{Cognitive Capability Captured} & \textbf{Observable Behavioral Pattern} & \textbf{Interpretive Logic} \\
\midrule
\textbf{GSB} (Goal--Strategy--Bug) &
Macro-level task understanding, planning, and execution &
Linear vs.\ cyclic trajectories &
Indicates whether the model maintains goal alignment and structured completion over time. \\
\addlinespace[2pt]
\textbf{Trajectory Archetype} &
Reasoning architecture prototype &
\emph{Linear Planner} / \emph{Iterative Validator} / \emph{Reactive Fixer} &
Contrasts direct plan–execute with iterative self-correction styles. \\
\addlinespace[2pt]
\textbf{Emergent Reflexivity (SCR, RD)} &
Self-monitoring and corrective reflex &
Frequency of ``wait'', ``check'', ``verify''-type utterances &
Signals meta-reasoning and self-diagnostic awareness. \\
\addlinespace[2pt]
\textbf{SDS} (Strategic Density Score) &
Density of key reasoning steps &
Effective reasoning tokens / total reasoning tokens &
Measures reasoning compactness and signal-to-noise ratio. \\
\addlinespace[2pt]
\textbf{TCR} (Trajectory Compression Ratio) &
Compressibility of reasoning path &
Accuracy retained after textual compression &
Evaluates redundancy vs.\ conciseness in the trajectory. \\
\bottomrule
\end{tabularx}
\caption{Mapping between behavioral metrics, cognitive capabilities, and interpretive logic. Weighted variants ($SDS_w$, $SCR_w$, $RD_w$) emphasize persistence and late-phase reflection.}
\label{tab:metric-interpretation}
\end{table}

\subsection{Model-wise Comparison and $\Delta M$ Index}

To compare two models under the same task, we define the composite performance
difference $\Delta M$ as a weighted sum of metric deltas:

\[
\Delta M =
0.35(\mathrm{SDS}_A - \mathrm{SDS}_B)
+ 0.35(\mathrm{SCR}_A - \mathrm{SCR}_B)
+ 0.20(E_A - E_B)
- 0.10(\mathrm{TCR}_A - \mathrm{TCR}_B).
\]

Here, SDS and SCR represent structured and verification densities,
$E$ is the emergent reasoning score,
and TCR (compression ratio) is subtracted to penalize overly terse responses.
The weight configuration $(0.35, 0.35, 0.20, 0.10)$
was chosen empirically to balance structural integrity (SDS, SCR)
and overall reasoning coherence (E)
while keeping linguistic brevity (TCR) as a secondary factor.

We classify model-level outcomes as:
\[
\text{Good: } \Delta M > 0.02, \quad
\text{Same: } |\Delta M| \le 0.02, \quad
\text{Bad: } \Delta M < -0.02.
\]

\section{Data Preprocessing and Sentence-Level Framework}
\label{appendix:data-preprocessing}

Raw conversation trajectories are stored as JSON arrays of events (\texttt{role}, \texttt{type}, timestamps, nested \texttt{message}).
We standardize roles (\texttt{assistant}, \texttt{user}, \texttt{system}); timestamps are used to compute total duration and assistant-active time.

\subsection{Sentence Extraction}
Let $\sentset(x)$ denote all sentences in text $x$. For a regex set $\Pi$:
\[
\density(x,\Pi) = \frac{\bigl|\{\,s\in\sentset(x): s \text{ matches } \Pi\,\}\bigr|}{|\sentset(x)|}, 
\qquad
\density(x,\Pi)=0 \text{ if } |\sentset(x)|=0.
\]

\subsection{Time Weighting}
Given assistant events $i$ with time-scale weight $w_i$, we define the time-weighted density:
\[
\densityw = \frac{\sum_i w_i \cdot \text{hits}_i}{\sum_i w_i \cdot \text{sentences}_i}.
\]
These preprocessing steps produce the normalized sentence-level corpora
(\texttt{ccbench\_parsed.csv} and \texttt{ccbench\_parsed\_weighted.csv}) used in all subsequent analyses.

\bibliographystyle{dcu}
% \bibliography{refs}

\end{document}
